% ==============================================================================
% Fichier : 09_evaluation_finale.tex
% Description : Évaluation Finale du cours de Hacking Éthique Avancé
% ==============================================================================

\chapter{Évaluation Finale du Cours}

Cette évaluation valide l'acquisition des compétences avancées en hacking éthique, couvrant les aspects juridiques, méthodologiques, techniques (recon, exploitation, post-exploitation) et professionnels (reporting).

% ------------------------------------------------------------------------------
\section{Partie A : QCM (Choix Multiples)}
% ------------------------------------------------------------------------------

\textbf{Consigne :} Cochez la ou les bonnes réponses.

\subsection*{Module 1 : Cadre Juridique et Méthodologique}
\begin{enumerate}
    \item En France, le cadre juridique principal encadrant les infractions informatiques est :
    \begin{itemize}
        \item[A.] Le RGPD
        \item[B.] Le Code pénal (Articles 323-1 et suivants)
        \item[C.] La loi Godfrain
        \item[D.] La norme ISO 27001
    \end{itemize}

    \item Quelle méthodologie divise le pentest en 7 phases distinctes ?
    \begin{itemize}
        \item[A.] OSSTMM
        \item[B.] PTES (Penetration Testing Execution Standard)
        \item[C.] NIST SP 800-115
        \item[D.] OWASP Testing Guide
    \end{itemize}

    \item Un test "Gray Box" signifie que :
    \begin{itemize}
        \item[A.] L'auditeur n'a aucune information préalable.
        \item[B.] L'auditeur dispose d'informations partielles (ex: compte utilisateur standard).
        \item[C.] L'auditeur a un accès complet (code source, configuration).
        \item[D.] L'auditeur ne dispose que d'une adresse IP publique.
    \end{itemize}

    \item La règle des "Règles d'Engagement" (ROE) définit principalement :
    \begin{itemize}
        \item[A.] Le prix de la prestation
        \item[B.] Les horaires et les actions interdites
        \item[C.] La méthode de nettoyage des traces
        \item[D.] La marque de café préférée du client
    \end{itemize}
\end{enumerate}

\subsection*{Module 2 : Reconnaissance et OSINT}
\begin{enumerate}
    \setcounter{enumi}{4}
    \item L'OSINT (Open Source Intelligence) se caractérise par :
    \begin{itemize}
        \item[A.] Une interaction directe avec les systèmes de la cible.
        \item[B.] La collecte d'informations publiques sans contact direct.
        \item[C.] L'utilisation exclusive d'outils de hacking.
        \item[D.] Un scan de ports intrusif.
    \end{itemize}

    \item Quel outil est principalement utilisé pour recher des vulnérabilités sur des documents publics via Google ?
    \begin{itemize}
        \item[A.] Shodan
        \item[B.] theHarvester
        \item[C.] Google Dorks
        \item[D.] Maltego
    \end{itemize}

    \item L'outil \texttt{Shodan} est utilisé pour :
    \begin{itemize}
        \item[A.] Scanner les ports locaux.
        \item[B.] Découvrir les appareils connectés à Internet.
        \item[C.] Cracker les mots de passe.
        \item[D.] Analyser le code source.
    \end{itemize}

    \item La commande \texttt{dig axfr} sert à :
    \begin{itemize}
        \item[A.] Effectuer un transfert de zone DNS (Zone Transfer).
        \item[B.] Scanner les ports TCP.
        \item[C.] Faire du brute force.
        \item[D.] Analyser les métadonnées d'une image.
    \end{itemize}
\end{enumerate}

\subsection*{Module 3 : Scanning et Énumération}
\begin{enumerate}
    \setcounter{enumi}{8}
    \item Le scan "SYN Scan" (-sS) est qualifié de :
    \begin{itemize}
        \item[A.] Complet et bruyant (logs).
        \item[B.] Furtif (demi-ouvert) et rapide.
        \item[C.] Impossible à réaliser sous Linux.
        \item[D.] Uniquement pour les protocoles UDP.
    \end{itemize}

    \item Dans Nmap, l'état "Filtered" signifie que :
    \begin{itemize}
        \item[A.] Le port est ouvert et une application écoute.
        \item[B.] Le port est fermé.
        \item[C.] Un pare-feu bloque l'accès.
        \item[D.] L'hôte est éteint.
    \end{itemize}

    \item L'énumération SMB (ports 139/445) permet de découvrir :
    \begin{itemize}
        \item[A.] Des partages réseaux et des utilisateurs.
        \item[B.] Des mots de passe en clair.
        \item[C.] La version exacte du BIOS.
        \item[D]. Les clés SSH.
    \end{itemize}

    \item L'outil \texttt{Dirb} est utilisé pour :
    \begin{itemize}
        \item[A.] Lancer des attaques DoS.
        \item[B]. Découvrir des répertoires et fichiers cachés sur un serveur web.
        \item[C.] Analyser le trafic réseau.
        \item[D.] Bruteforcer des identifiants SSH.
    \end{itemize}
\end{enumerate}

\subsection*{Module 4 : Exploitation}
\begin{enumerate}
    \setcounter{enumi}{12}
    \item Une injection SQL de type "Blind" signifie que :
    \begin{itemize}
        \item[A.] L'attaquant peut voir les données directement sur la page.
        \item[B.] L'attaquant ne voit pas les données mais peut les déduire via des messages d'erreur ou le temps de réponse.
        \item[C.] L'attaquant peut exécuter des commandes système.
        \item[D.] L'attaquant peut écrire dans la base de données mais pas lire.
    \end{itemize}

    \item Quel outil est le standard pour l'exploitation de vulnérabilités ?
    \begin{itemize}
        \item[A.] Nmap
        \item[B.] Metasploit Framework
        \item[C.] Wireshark
        \item[D.] Burp Suite
    \end{itemize}

    \item Un "Reverse Shell" fonctionne de la manière suivante :
    \begin{itemize}
        \item[A.] L'attaquant se connecte à la cible sur un port spécifique.
        \item[B.] La cible se connecte à l'attaquant sur un port spécifique.
        \item[C.] L'attaquant envoie un fichier à la cible.
        \item[D.] La cible télécharge une mise à jour.
    \end{itemize}

    \item L'outil \texttt{SQLMap} est utilisé pour automatiser :
    \begin{itemize}
        \item[A.] Les scans de ports.
        \item[B.] Les injections SQL.
        \item[C.] Les attaques XSS.
        \item[D.] L'énumération DNS.
    \end{itemize}

    \item Quelle est la différence entre LFI et RFI ?
    \begin{itemize}
        \item[A.] LFI est à distance, RFI est local.
        \item[B.] LFI inclut des fichiers locaux, RFI inclut des fichiers distants.
        \item[C.] LFI est pour Linux, RFI est pour Windows.
        \item[D.] Il n'y a pas de différence.
    \end{itemize}
\end{enumerate}

\subsection*{Module 5 : Post-Exploitation et Mouvement Latéral}
\begin{enumerate}
    \setcounter{enumi}{17}
    \item L'escalade de privilèges consiste à :
    \begin{itemize}
        \item[A.] Passer d'un compte standard à un compte administrateur/root.
        \item[B.] Passer d'un serveur à un autre.
        \item[C.] Changer le mot de passe d'un utilisateur.
        \item[D.] Effacer les logs.
    \end{itemize}

    \item L'outil \texttt{Mimikatz} est célèbre pour :
    \begin{itemize}
        \item[A.] Scanner les vulnérabilités web.
        \item[B.] Extraire les mots de passe et hashs de la mémoire Windows (LSASS).
        \item[C.] Effectuer des scans de ports furtifs.
        \item[D.] Analyser les fichiers PDF.
    \end{itemize}

    \item La technique "Pass-the-Hash" permet de :
    \begin{itemize}
        \item[A.] Craquer un hash pour trouver le mot de passe en clair.
        \item[B.] S'authentifier en utilisant directement le hash du mot de passe.
        \item[C.] Modifier le hash dans la base SAM.
        \item[D.] Créer un nouvel utilisateur.
    \end{itemize}

    \item L'outil \texttt{BloodHound} sert à :
    \begin{itemize}
        \item[A.] Surveiller le trafic réseau en temps réel.
        \item[B.] Visualiser les relations d'attaque dans un environnement Active Directory.
        \item[C.] Détecter les malwares sur le disque.
        \item[D.] Renforcer les pare-feu.
    \end{itemize}

    \item Le "Pivoting" consiste à :
    \begin{itemize}
        \item[A.] Tourner le disque dur pour un accès physique.
        \item[B.] Utiliser une machine compromise pour attaquer un réseau inaccessible.
        \item[C.] Retourner à la case de l'attaquant.
        \item[D.] Changer d'adresse IP publique.
    \end{itemize}
\end{enumerate}

\subsection*{Module 6 : Évasion et Anti-Forensics}
\begin{enumerate}
    \setcounter{enumi}{22}
    \item L'antivirus détecte généralement Mimikatz via :
    \begin{itemize}
        \item[A.] La détection comportementale.
        \item[B.] La signature de fichier (statique).
        \item[C.] L'analyse du disque dur exclusivement.
        \item[D.] L'analyse du cache navigateur.
    \end{itemize}

    \item La technique "Living Off The Land" (LoL) implique d'utiliser :
    \begin{itemize}
        \item[A.] Des outils tiers téléchargés depuis GitHub.
        \item[B.] Des outils légitimes du système (ex: Certutil, PowerShell) pour l'attaque.
        \item[C.] Des outils payants.
        \item[D.] Un live CD Linux.
    \end{itemize}

    \item Le but de l'obfuscation de code est de :
    \begin{itemize}
        \item[A.] Améliorer la lisibilité du code.
        \item[B.] Rendre le code plus rapide.
        \item[C.] Modifier le code pour éviter la détection par signature.
        \item[D.] Compresser le code.
    \end{itemize}

    \item AMSI (Antimalware Scan Interface) est conçu pour :
    \begin{itemize}
        \item[A.] Bloquer les ports non utilisés.
        \item[B.] Scanner le contenu des scripts (PowerShell, VBS) avant exécution.
        \item[C.] Filtrer les emails.
        \item[D.] Chiffrer le disque dur.
    \end{itemize}
\end{enumerate}

\subsection*{Module 7 : Rapport et Remédiation}
\begin{enumerate}
    \setcounter{enumi}{26}
    \item La partie "Executive Summary" d'un rapport est destinée à :
    \begin{itemize}
        \item[A.] Les équipes techniques.
        \item[B.] Les développeurs.
        \item[C.] La direction (PDG, RSSI).
        \item[D.] Le support client.
    \end{itemize}

    \item Le score CVSS mesure :
    \begin{itemize}
        \item[A.] Le coût financier d'une faille.
        \item[B.] La gravité d'une vulnérabilité (de 0 à 10).
        \item[C.] Le temps de correction.
        \item[D.] La complexité du code.
    \end{itemize}

    \item Quelle est la meilleure pratique pour proposer une correction ?
    \begin{itemize}
        \item[A.] Corriger la faille immédiatement sans dire au client.
        \item[B.] Proposer une compensation (ex: bloquer le port) si la mise à jour n'est pas possible.
        \item[C.] Ignorer la faille si elle est difficile à exploiter.
        \item[D.] Demander une augmentation de budget.
    \end{itemize}

    \item La preuve de concept (PoC) dans un rapport sert à :
    \begin{itemize}
        \item[A.] Prouver que l'attaquant n'a pas menti sur la compromission.
        \item[B.] Montrer le code source de l'outil utilisé.
        \item[C.] Faire de la publicité pour un outil.
        \item[D.] Remplir des pages.
    \end{itemize}
\end{enumerate}

% ------------------------------------------------------------------------------
\section{Partie B : Vrai ou Faux}
% ------------------------------------------------------------------------------

\begin{enumerate}
    \item \textbf{V/F} : Il est légal de scanner n'importe quel site web sans autorisation. \textbf{(Faux - Le scan actif nécessite une autorisation préalable)}
    \item \textbf{V/F} : Le "SYN Scan" est plus furtif que le "Connect Scan". \textbf{(Vrai - Il n'établit pas la connexion complète)}
    \item \textbf{V/F} : Une attaque "Black Box" signifie que l'attaquant a accès au code source. \textbf{(Faux - C'est "White Box")}
    \item \textbf{V/F} : Le Pass-the-Hash fonctionne principalement avec le protocole Kerberos. \textbf{(Faux - C'est NTLM)}
    \item \textbf{V/F} : L'injection SQL de type "Error Based" affiche des messages d'erreur SQL sur la page web. \textbf{(Vrai)}
    \item \textbf{V/F} : L'outil \texttt{Hydra} est utilisé pour les injections SQL. \textbf{(Faux - Il sert au brute forcing de mots de passe)}
    \item \textbf{V/F} : Le "Kernel Exploit" est une technique d'escalade de privilèges Linux. \textbf{(Vrai)}
    \item \textbf{V/F} : Effacer les logs Windows (\texttt{wevtutil cl}) est une action silencieuse et indétectable. \textbf{(Faux - L'effacement crée un événement "Log Cleared")}
    \item \textbf{V/F} : Le "Pivoting" permet d'accéder à des réseaux internes via une machine compromise en DMZ. \textbf{(Vrai)}
    \item \textbf{V/F} : Un score CVSS de 2.0 indique une vulnérabilité critique. \textbf{(Faux - C'est faible. Critique est généralement > 9.0)}
\end{enumerate}

% ------------------------------------------------------------------------------
\section{Partie C : Questions Ouvertes}
% ------------------------------------------------------------------------------

\begin{enumerate}
    \item Expliquez la différence fondamentale entre un "Reverse Shell" et un "Bind Shell" en termes de connexion réseau et contournement de pare-feu.
    \item Décrivez les 3 étapes de l'exploitation d'une injection SQL de type "Union Based" (Détection, Nombre de colonnes, Type de données).
    \item Pourquoi l'utilisation d'outils légitimes ("Living Off The Land") est-elle plus difficile à détecter que le dépôt d'un exécutable externe (comme Mimikatz.exe) ?
    \item Expliquez le concept de "Pivoting" et pourquoi il est nécessaire lors d'un pentest d'une architecture avec une DMZ et un LAN.
    \item Quelle est la différence entre la détection statique (signature) et la détection comportementale (heuristique) dans un antivirus ?
    \item Pourquoi est-il recommandé de ne pas éteindre un serveur compromis lors d'une investigation forensique ?
    \item Expliquez ce qu'est le "Null Session" en énumération SMB et quel type d'information il permet de récupérer.
    \item Définissez le concept de "Credential Dumping" et citez un outil majeur pour le réaliser sous Windows.
\end{enumerate}

% ------------------------------------------------------------------------------
\section{Partie D : Cas Pratique - Capture The Flag}
% ------------------------------------------------------------------------------

\subsection{Scénario : Audit du Serveur Web de LiZbiZ}

\begin{displayquote}
Vous êtes connectés au réseau interne de LiZbiZ (Lab GNS3). On vous demande de compromettre le serveur Web interne (IP : 192.168.1.50) et de récupérer le fichier confidentiel stocké sur le serveur de fichiers interne.
\end{displayquote}

\textbf{Questions et tâches :}

\begin{enumerate}
    \item \textbf{Reconnaissance :} Utilisez \texttt{nmap} pour identifier le système d'exploitation et les services ouverts sur la cible (192.168.1.50). Quels sont les ports ouverts ?
    \item \textbf{Énumération :} En utilisant un navigateur ou \texttt{curl}, identifiez le CMS utilisé (ex: WordPress). Quels fichiers/dossiers sensibles sont accessibles (Directory Busting) ?
    \item \textbf{Exploitation :} Vous découvrez une page de login vulnérable. Proposez une attaque par injection SQL (ou brute force si simple) pour accéder au compte administrateur.
    \item \textbf{Post-Exploitation :} Une fois connecté, comment obtenez-vous un shell sur le serveur ? (Indice : Web Shell ou exécution de commande).
    \item \textbf{Élévation de Privilèges :} Vérifiez les permissions sudo ou les capacités du kernel. Comment passez-vous de l'utilisateur 'www-data' à 'root' ?
    \item \textbf{Preuve :} Récupérez le contenu du fichier \texttt{/root/flag.txt}. Indiquez le hash contenu.
\end{enumerate}